% GENERATED FILE. Edit canonical lessons, then run npm run generate:books.
% canonical-source-hash: fnv1a64:bb2f955f1d0181be
% canonical-lessons: TE-C426-sankocincu, TE-C426-siddhapadu, TE-C426-asrayincu, TE-C426-uddesapurvakanga, TE-C426-arambham

\chapter{Before You Act}
\label{ch:te-before-you-act}
\hlchaptermodality{426}

\begin{chapteropening}
\textbf{By the end of this chapter:} \emph{I can talk about hesitating, getting ready and acting on purpose.}

Complete the last lesson of chapter 426: I can talk about hesitating, getting ready and acting on purpose.
\end{chapteropening}

\section[\texorpdfstring{\te{సంకోచించు} (saṅkōciñcu)}{సంకోచించు (saṅkōciñcu)}]{\te{సంకోచించు} (saṅkōciñcu) — to hold back; to hesitate}

\label{lesson:TE-C426-sankocincu}



\begin{quote}
Before the new one: say the Telugu for to finalise, then the Telugu for to reconsider.
\end{quote}

\subsection*{You'll want to know: \te{సంకోచించు}}
\textbf{\te{సంకోచించు}} — \emph{saṅkōciñcu} — \textquotedblleft{}to hold back; to hesitate\textquotedblright{}.

\begin{quote}
\te{ఆమె} \te{నిజం} \te{చెప్పడానికి} \te{సంకోచించింది}.
\end{quote}

\begin{quote}
\emph{āme nijaṁ ceppaḍāniki saṅkōciñcindi.}
\end{quote}

\begin{quote}
She hesitated to tell the truth.
\end{quote}

Past (I …): \te{సంకోచించాను}, \emph{saṅkōciñcānu}.

What you hesitate over takes -\te{డానికి}. Close to \te{వెనకాడు}, which you know.

\begin{cousinweb}
\textbf{\te{సంకోచించు}} (\emph{saṅkōciñcu}) is built on Sanskrit \emph{saṁkoca}, \textquotedblleft{}contraction, shrinking; shyness\textquotedblright{} (Monier-Williams), with the Telugu verb ending -\te{ఇంచు}.
\end{cousinweb}

\begin{grammarlens}[title={What you've built}]
One more word for saying what happens if.
\end{grammarlens}

\subsection*{Your turn}
\begin{itemize}
\raggedright
  \item \emph{Say it:} \emph{saṅkōciñcu}
  \item \emph{Say it:} \emph{saṅkōciñcu}, once more
  \item \emph{Say it:} \emph{āme nijaṁ ceppaḍāniki saṅkōciñcindi.} — She hesitated to tell the truth.
  \item \emph{Recall:} say the Telugu for to finalise, then the Telugu for to reconsider, then say \emph{saṅkōciñcu} again
\end{itemize}

\begin{tcolorbox}[breakable,colback=teal!4,colframe=teal!35!black,title={Before you move on}]
What does \te{సంకోచించు} mean? (\textquotedblleft{}To hold back; to hesitate\textquotedblright{}.) Say it once more.
\end{tcolorbox}
\section[\texorpdfstring{\te{సిద్ధపడు} (siddhapaḍu)}{సిద్ధపడు (siddhapaḍu)}]{\te{సిద్ధపడు} (siddhapaḍu) — to be willing; to get ready}

\label{lesson:TE-C426-siddhapadu}



\begin{quote}
Before the new one: say the Telugu for to reconsider, then the Telugu for to hold back.
\end{quote}

\subsection*{You'll want to know: \te{సిద్ధపడు}}
\textbf{\te{సిద్ధపడు}} — \emph{siddhapaḍu} — \textquotedblleft{}to be willing; to get ready\textquotedblright{}.

\begin{quote}
\te{అతను} \te{ఎంత} \te{దూరమైనా} \te{వెళ్ళడానికి} \te{సిద్ధపడ్డాడు}.
\end{quote}

\begin{quote}
\emph{atanu enta dūramainā veḷḷaḍāniki siddhapaḍḍāḍu.}
\end{quote}

\begin{quote}
He was ready to go however far it took.
\end{quote}

Past (I …): \te{సిద్ధపడ్డాను}, \emph{siddhapaḍḍānu}.

\te{సిద్ధం}, ready, + \te{పడు}.

\begin{cousinweb}
\textbf{\te{సిద్ధపడు}} (\emph{siddhapaḍu}) is built from Sanskrit \emph{siddha}, \textquotedblleft{}accomplished; prepared\textquotedblright{} (CDIAL 13401), the word in \te{సిద్ధం}, and the native verb \emph{paḍu} (DEDR 3853).
\end{cousinweb}

\begin{grammarlens}[title={What you've built}]
One more word for saying what happens if.
\end{grammarlens}

\subsection*{Your turn}
\begin{itemize}
\raggedright
  \item \emph{Say it:} \emph{siddhapaḍu}
  \item \emph{Say it:} \emph{siddhapaḍu}, once more
  \item \emph{Say it:} \emph{atanu enta dūramainā veḷḷaḍāniki siddhapaḍḍāḍu.} — He was ready to go however far it took.
  \item \emph{Recall:} say the Telugu for to reconsider, then the Telugu for to hold back, then say \emph{siddhapaḍu} again
\end{itemize}

\begin{tcolorbox}[breakable,colback=teal!4,colframe=teal!35!black,title={Before you move on}]
What does \te{సిద్ధపడు} mean? (\textquotedblleft{}To be willing; to get ready\textquotedblright{}.) Say it once more.
\end{tcolorbox}
\section[\texorpdfstring{\te{ఆశ్రయించు} (āśrayiñcu)}{ఆశ్రయించు (āśrayiñcu)}]{\te{ఆశ్రయించు} (āśrayiñcu) — to turn to; to resort to}

\label{lesson:TE-C426-asrayincu}



\begin{quote}
Before the new one: say the Telugu for to hold back, then the Telugu for to be willing.
\end{quote}

\subsection*{You'll want to know: \te{ఆశ్రయించు}}
\textbf{\te{ఆశ్రయించు}} — \emph{āśrayiñcu} — \textquotedblleft{}to turn to; to resort to\textquotedblright{}.

\begin{quote}
\te{న్యాయం} \te{కోసం} \te{వాళ్ళు} \te{కోర్టును} \te{ఆశ్రయించారు}.
\end{quote}

\begin{quote}
\emph{nyāyaṁ kōsaṁ vāḷḷu kōrṭunu āśrayiñcāru.}
\end{quote}

\begin{quote}
They turned to the court for justice.
\end{quote}

Past (I …): \te{ఆశ్రయించాను}, \emph{āśrayiñcānu}.

For turning to someone for help or shelter.

\begin{cousinweb}
\textbf{\te{ఆశ్రయించు}} (\emph{āśrayiñcu}) is built on Sanskrit \emph{āśraya}, \textquotedblleft{}resting place, refuge\textquotedblright{} (CDIAL 1466), with the Telugu verb ending -\te{ఇంచు}; \te{విమానాశ్రయం}, an airport, ends in the same word.
\end{cousinweb}

\begin{grammarlens}[title={What you've built}]
One more word for saying what happens if.
\end{grammarlens}

\subsection*{Your turn}
\begin{itemize}
\raggedright
  \item \emph{Say it:} \emph{āśrayiñcu}
  \item \emph{Say it:} \emph{āśrayiñcu}, once more
  \item \emph{Say it:} \emph{nyāyaṁ kōsaṁ vāḷḷu kōrṭunu āśrayiñcāru.} — They turned to the court for justice.
  \item \emph{Recall:} say the Telugu for to hold back, then the Telugu for to be willing, then say \emph{āśrayiñcu} again
\end{itemize}

\begin{tcolorbox}[breakable,colback=teal!4,colframe=teal!35!black,title={Before you move on}]
What does \te{ఆశ్రయించు} mean? (\textquotedblleft{}To turn to; to resort to\textquotedblright{}.) Say it once more.
\end{tcolorbox}
\section[\texorpdfstring{\te{ఉద్దేశపూర్వకంగా} (uddēśapūrvakaṅgā)}{ఉద్దేశపూర్వకంగా (uddēśapūrvakaṅgā)}]{\te{ఉద్దేశపూర్వకంగా} (uddēśapūrvakaṅgā) — deliberately; on purpose}

\label{lesson:TE-C426-uddesapurvakanga}



\begin{quote}
Before the new one: say the Telugu for to be willing, then the Telugu for to turn to.
\end{quote}

\subsection*{You'll want to know: \te{ఉద్దేశపూర్వకంగా}}
\textbf{\te{ఉద్దేశపూర్వకంగా}} — \emph{uddēśapūrvakaṅgā} — \textquotedblleft{}deliberately; on purpose\textquotedblright{}.

\begin{quote}
\te{అతను} \te{ఉద్దేశపూర్వకంగా} \te{ఆ} \te{విషయం} \te{దాచాడు}.
\end{quote}

\begin{quote}
\emph{atanu uddēśapūrvakaṅgā ā viṣayaṁ dācāḍu.}
\end{quote}

\begin{quote}
He hid the matter on purpose.
\end{quote}

\te{ఉద్దేశం}, an intention, + \te{పూర్వకంగా}, with ... in front.

\begin{cousinweb}
\textbf{\te{ఉద్దేశపూర్వకంగా}} (\emph{uddēśapūrvakaṅgā}) is built from Sanskrit \emph{uddeśa}, \textquotedblleft{}aim, direction\textquotedblright{}, and \emph{pūrvaka}, \textquotedblleft{}earlier, preceded by\textquotedblright{} (Monier-Williams), with the adverb ending -\te{గా}: with the aim already in front.
\end{cousinweb}

\begin{grammarlens}[title={What you've built}]
One more word for saying what happens if.
\end{grammarlens}

\subsection*{Your turn}
\begin{itemize}
\raggedright
  \item \emph{Say it:} \emph{uddēśapūrvakaṅgā}
  \item \emph{Say it:} \emph{uddēśapūrvakaṅgā}, once more
  \item \emph{Say it:} \emph{atanu uddēśapūrvakaṅgā ā viṣayaṁ dācāḍu.} — He hid the matter on purpose.
  \item \emph{Recall:} say the Telugu for to be willing, then the Telugu for to turn to, then say \emph{uddēśapūrvakaṅgā} again
\end{itemize}

\begin{tcolorbox}[breakable,colback=teal!4,colframe=teal!35!black,title={Before you move on}]
What does \te{ఉద్దేశపూర్వకంగా} mean? (\textquotedblleft{}Deliberately; on purpose\textquotedblright{}.) Say it once more.
\end{tcolorbox}
\section[\texorpdfstring{\te{ఆరంభం} (ārambhaṁ)}{ఆరంభం (ārambhaṁ)}]{\te{ఆరంభం} (ārambhaṁ) — a beginning; a start}

\label{lesson:TE-C426-arambham}



\begin{quote}
Before the new one: say the Telugu for to turn to, then the Telugu for deliberately.
\end{quote}

\subsection*{You'll want to know: \te{ఆరంభం}}
\textbf{\te{ఆరంభం}} — \emph{ārambhaṁ} — \textquotedblleft{}a beginning; a start\textquotedblright{}.

\begin{quote}
\te{ఆరంభంలో} \te{పని} \te{కష్టంగా} \te{ఉంది}.
\end{quote}

\begin{quote}
\emph{ārambhaṁlō pani kaṣṭaṅgā undi.}
\end{quote}

\begin{quote}
At the start the work is hard.
\end{quote}

Formal beside \te{మొదలు}; \te{ఆరంభించు}, to begin.

\begin{cousinweb}
\textbf{\te{ఆరంభం}} (\emph{ārambhaṁ}) is borrowed from Sanskrit \emph{ārambha}, \textquotedblleft{}beginning\textquotedblright{} (CDIAL 1307).
\end{cousinweb}

\begin{grammarlens}[title={What you've built}]
One more word for saying what happens if.
\end{grammarlens}

\subsection*{Your turn}
\begin{itemize}
\raggedright
  \item \emph{Say it:} \emph{ārambhaṁ}
  \item \emph{Say it:} \emph{ārambhaṁ}, once more
  \item \emph{Say it:} \emph{ārambhaṁlō pani kaṣṭaṅgā undi.} — At the start the work is hard.
  \item \emph{Recall:} say the Telugu for to turn to, then the Telugu for deliberately, then say \emph{ārambhaṁ} again
\end{itemize}

\begin{tcolorbox}[breakable,colback=teal!4,colframe=teal!35!black,title={Before you move on}]
What does \te{ఆరంభం} mean? (\textquotedblleft{}A beginning; a start\textquotedblright{}.) Say it once more.
\end{tcolorbox}
